\newpage
\section{Common Failure Cases of Descriptive Questions}
\label{sec:desc_failure}

We provide 30 concrete examples on each of the descriptive questions in which the vast majority of representative models fail to provide the correct answer. Common failures of models include:
\vspace{-1ex}
\begin{itemize}
    \item Models cannot correctly \textbf{localize} subplot when \textbf{many subplots are present} (\cref{subsec:desc_example1,,subsec:desc_example10,,subsec:desc_example11,,subsec:desc_example18,,subsec:desc_example26}).
    \item Models use \textbf{incorrect elements} of the charts to provide an answer (\cref{subsec:desc_example2,,subsec:desc_example4,,subsec:desc_example5,,subsec:desc_example6,,subsec:desc_example7,,subsec:desc_example8,,subsec:desc_example9,,subsec:desc_example10,,subsec:desc_example12,,subsec:desc_example15,,subsec:desc_example22,,subsec:desc_example25}).
    \item Models make \textbf{OCR mistakes} (\cref{subsec:desc_example2,,subsec:desc_example3,,subsec:desc_example5,,subsec:desc_example9,,subsec:desc_example15,,subsec:desc_example19}).
    \item Models fail to identify relevant elements when they are \textbf{not} close to the subplot (\cref{subsec:desc_example3,,subsec:desc_example23}).
    \item Models \textbf{hallucinate} (\cref{subsec:desc_example13,,subsec:desc_example16,,subsec:desc_example17,,subsec:desc_example20,,subsec:desc_example24,,subsec:desc_example26}).
    \item Models fail to tackle \textbf{tricky or unconventional scenarios} (\cref{subsec:desc_example14,,subsec:desc_example21,,subsec:desc_example27}).
    \item Models fail to \textbf{count} (\cref{subsec:desc_example15,,subsec:desc_example28,,subsec:desc_example29,,subsec:desc_example30}).
\end{itemize}

\begin{table*}[h]
  \centering
  \caption{Overview of failure case examples in descriptive questions. We provide 30 examples for each of the descriptive questions with both answerable and unanswerable scenarios.}
  \scalebox{0.9}{
  \begin{tabular}{ccc@{}}
    \hline
    \toprule

    \textbf{Example} & \textbf{QID} & \textbf{Answerable}  \\
    \midrule
1 & \multirow{2}{*}{1} & \cmark   \\
2 &  & \xmark\\
\midrule
3 & \multirow{2}{*}{2} & \cmark   \\
4 &  & \xmark \\
\midrule
5 & \multirow{2}{*}{3} & \cmark   \\
6 &  & \xmark  \\
\midrule
7 & 4 & \cmark   \\
\midrule
8 & 5 & \cmark   \\
\midrule
9 & 6 & \cmark   \\
\midrule
10 & 7 & \cmark   \\
\midrule
11 & \multirow{2}{*}{8} & \cmark   \\
12 &  & \xmark  \\
\midrule
13 & \multirow{2}{*}{9} & \cmark   \\
14 &  & \xmark  \\
\midrule
15 & \multirow{2}{*}{10} & \cmark   \\
16 &  & \xmark  \\
\midrule
17 & \multirow{2}{*}{11} & \cmark   \\
18 &  & \xmark  \\
\midrule
19 & \multirow{2}{*}{12} & \cmark   \\
20 &  & \xmark  \\
\midrule
21 & \multirow{2}{*}{13} & \cmark   \\
22 &  & \xmark \\
\midrule
23 & \multirow{2}{*}{14} & \cmark   \\
24 &  & \xmark  \\
\midrule
25 & \multirow{2}{*}{15} & \cmark   \\
26 &  & \xmark  \\
\midrule
27 & 16 & \cmark   \\
\midrule
28 & 17 & \cmark   \\
\midrule
29 & 18 & \cmark   \\
\midrule
30 & 19 & \cmark   \\
    \bottomrule
    \hline
    
  \end{tabular}
  }
  \label{tab:desc_failure_cases}
\end{table*}
\newpage
\subsection{Example 1}
\label{subsec:desc_example1}
\begin{example}[Example]
\includegraphics[width=0.5\linewidth]{appendix_desc/159.jpg}
\begin{lstlisting}
Question: For the subplot at row 4 and column 4, what is its title?
    * Your final answer should be the most relevant title of the plot that is explicitly written.
    * If the plot does not have an explicit title or contains only a letter, answer 'Not Applicable'.
    
Answer: Sensitivity
\end{lstlisting}
\end{example}

\begin{wronganswer}
\textbf{GPT-4o}: \small{The title of the subplot at row 4 and column 4 is \red{"Sadness"}.}
\end{wronganswer}

\begin{wronganswer}
\textbf{Claude 3 Sonnet}: \small{The title of the subplot at row 4 and column 4 is \red{Fear}.}
\end{wronganswer}

\begin{wronganswer}
\textbf{Reka Core}: \small{ The title of the subplot at row 4 and column 4 is \red{"Pain"}. This can be determined by looking at the label at the top of the subplot, which is a common practice in scientific figures to indicate the category or variable being represented in each individual plot.}
\end{wronganswer}

\begin{wronganswer}
\textbf{InternVL Chat V1.5}: \small{\red{Fear}}
\end{wronganswer}

\begin{wronganswer}
\textbf{MGM HD Yi 34B}: \small{The subplot at row 4 and column 4 is titled \red{"Pain"}. This title is explicitly written at the top of the subplot, indicating the emotional category being represented by the scatter plot in that particular position.}
\end{wronganswer}

\begin{wronganswer}
\textbf{IDEFICS 2}: \small{\red{Fatigue}.}
\end{wronganswer}
\newpage

\subsection{Example 2}
\label{subsec:desc_example2}
\begin{example}[Example]
\includegraphics[width=0.5\linewidth]{appendix_desc/1822.jpg}
\begin{lstlisting}
Question: For the subplot at row 3 and column 1, what is its title?
    * Your final answer should be the most relevant title of the plot that is explicitly written.
    * If the plot does not have an explicit title or contains only a letter, answer 'Not Applicable'.
    
Answer: Not Applicable
\end{lstlisting}
\end{example}

\begin{wronganswer}
\textbf{GPT-4o}: \small{\red{Powell Q\&A}}
\end{wronganswer}

\begin{wronganswer}
\textbf{Claude 3 Sonnet}: \small{The title of the subplot in the third row is \red{"Avg Price (\$/Bbl)"}.}
\end{wronganswer}

\begin{wronganswer}
\textbf{Reka Core}: \small{ \red{Avg. Price \$ [157.0 159.0]}.}
\end{wronganswer}

\begin{wronganswer}
\textbf{InternVL Chat V1.5}: \small{\red{FOMC ann}.}
\end{wronganswer}

\begin{wronganswer}
\textbf{MGM HD Yi 34B}: \small{The title of the subplot at row 3 and column 1 is \red{"Powell Q\&A"}. This title is explicitly written in the upper right corner of the subplot, indicating the subject or event that the data is related to.}
\end{wronganswer}

\begin{wronganswer}
\textbf{IDEFICS 2}: \small{\red{Powell Q\&A}.}
\end{wronganswer}
\newpage
\subsection{Example 3}
\label{subsec:desc_example3}
\begin{example}[Example]
\includegraphics[width=0.5\linewidth]{appendix_desc/1241.jpg}
\begin{lstlisting}
Question: For the subplot at row 2 and column 1, what is the label of the x-axis?
    * Your final answer should be the label of the x-axis that is explicitly written, including the case when x-axis is shared across multiple subplots. When the x-axis is present on both the top and bottom of the plot, answer the label of the x-axis at the bottom.
    * If the plot does not have an explicit x-axis label, answer 'Not Applicable'.
    
Answer: BJD-2450000
\end{lstlisting}
\end{example}

\begin{wronganswer}
\textbf{GPT-4o}: \small{The label of the x-axis for the subplot at row 2 and column 1 is \red{"Not Applicable"}.}
\end{wronganswer}

\begin{wronganswer}
\textbf{Claude 3 Sonnet}: \small{The x-axis label for the subplot at row 2, column 1 is \red{Not Applicable. There is no explicit x-axis label written for that particular subplot.}}
\end{wronganswer}

\begin{wronganswer}
\textbf{Reka Core}: \small{ The label of the x-axis for the subplot at row 2 and column 1 is \red{'H$\alpha$'}. This is \red{explicitly written at the bottom of the subplot, indicating that the x-axis represents the H-alpha spectral line, which is a specific deep red visible spectral line created by hydrogen.}}
\end{wronganswer}

\begin{wronganswer}
\textbf{InternVL Chat V1.5}: \small{The label of the x-axis for the subplot at row 2 and column 1 is \red{"Ha"}.}
\end{wronganswer}

\begin{wronganswer}
\textbf{MGM HD Yi 34B}: \small{The label of the x-axis in the subplot at row 2 and column 1 is \red{'BJD-245000'}. This label is explicitly written at the bottom of the plot, indicating the time scale used for the data points in this subplot.}
\end{wronganswer}

\begin{wronganswer}
\textbf{IDEFICS 2}: \small{\red{BJD-245000}.}
\end{wronganswer}
\newpage
\subsection{Example 4}
\label{subsec:desc_example4}
\begin{example}[Example]
\includegraphics[width=0.5\linewidth]{appendix_desc/493.jpg}
\begin{lstlisting}
Question: For the subplot at row 2 and column 1, what is the label of the x-axis?
    * Your final answer should be the label of the x-axis that is explicitly written, including the case when x-axis is shared across multiple subplots. When the x-axis is present on both the top and bottom of the plot, answer the label of the x-axis at the bottom.
    * If the plot does not have an explicit x-axis label, answer 'Not Applicable'.
    
Answer: Not Applicable
\end{lstlisting}
\end{example}

\begin{wronganswer}
\textbf{GPT-4o}: \small{The label of the x-axis for the subplot at row 2 and column 1 is "\red{U}".}
\end{wronganswer}

\begin{wronganswer}
\textbf{Claude 3 Sonnet}: \small{The label of the x-axis for the subplot at row 2, column 1 is\red{Stock}.}
\end{wronganswer}

\begin{wronganswer}
\textbf{Reka Core}: \small{ The label of the x-axis for the subplot at row 2 and column 1 is '\red{U}', which is explicitly written \red{on the bottom} of the plot. This label is \red{shared across multiple subplots in the same row, indicating that the x-axis represents a common variable or time period for these plots}.}
\end{wronganswer}

\begin{wronganswer}
\textbf{InternVL Chat V1.5}: \small{The label of the x-axis for the subplot at row 2 and column 1 is "\red{Stock}."}
\end{wronganswer}

\begin{wronganswer}
\textbf{MGM HD Yi 34B}: \small{The label of the x-axis in the subplot at row 2 and column 1 is "\red{Stock}". This label is explicitly written at the bottom of the plot, indicating the time series data points are measured against the stock variable.}
\end{wronganswer}

\begin{wronganswer}
\textbf{IDEFICS 2}: \small{\red{Stock}.}
\end{wronganswer}
\newpage
\subsection{Example 5}
\label{subsec:desc_example5}
\begin{example}[Example]
\includegraphics[width=0.5\linewidth]{appendix_desc/173.jpg}
\begin{lstlisting}
Question: For the subplot at row 1 and column 2, what is the label of the y-axis?
    * Your final answer should be the label of the y-axis that is explicitly written, including the case when y-axis is shared across multiple subplots. When the y-axis is present on both the left and right of the plot, answer the label of the y-axis at the left.
    * If the plot does not have an explicit y-axis label, answer 'Not Applicable'.
Answer: LLC (Ways)

Note: several models answered it incorrectly because they did not follow the instruction to focus on the label of the y-axis at the left (instead of the right) when labels are present at both the left and right.
\end{lstlisting}
\end{example}

\begin{wronganswer}
\textbf{GPT-4o}: \small{The label of the y-axis for the subplot at row 1 and column 2 is \red{"MBW (\%)"}.}
\end{wronganswer}

\begin{wronganswer}
\textbf{Claude 3 Sonnet}: \small{The y-axis label for the subplot at row 1 and column 2 is LLC \red{(MB/s)}.}
\end{wronganswer}

\begin{wronganswer}
\textbf{Reka Core}: \small{ The label of the y-axis for the subplot at row 1 and column 2 is '\red{MBW (\%)}', which stands for Memory Bandwidth Utilization percentage. This is \red{explicitly written on the left side} of the y-axis for this particular subplot.}
\end{wronganswer}

\begin{wronganswer}
\textbf{InternVL Chat V1.5}: \small{The label of the y-axis for the subplot at row 1 and column 2 is "\red{Load (Gbps)}."}
\end{wronganswer}

\begin{wronganswer}
\textbf{MGM HD Yi 34B}: \small{The label of the y-axis in the subplot at row 1 and column 2 is "\red{Freq (GHz)}". This label is explicitly written and is shared across all \red{three} subplots in the "Freq Allocation" section of the image.}
\end{wronganswer}

\begin{wronganswer}
\textbf{IDEFICS 2}: \small{\red{Load \& QoS}.}
\end{wronganswer}
\newpage
\subsection{Example 6}
\label{subsec:desc_example6}
\begin{example}[Example]
\includegraphics[width=0.5\linewidth]{appendix_desc/284.jpg}
\begin{lstlisting}
Question: For the subplot at row 1 and column 3, what is the label of the y-axis?
    * Your final answer should be the label of the y-axis that is explicitly written, including the case when y-axis is shared across multiple subplots. When the y-axis is present on both the left and right of the plot, answer the label of the y-axis at the left.
    * If the plot does not have an explicit y-axis label, answer 'Not Applicable'.
Answer: Not Applicable
\end{lstlisting}
\end{example}

\begin{correctanswer}
\textbf{GPT-4o}: \small{The label of the y-axis for the subplot at row 1 and column 3 is '\textbf{Not Applicable}'.}
\end{correctanswer}

\begin{wronganswer}
\textbf{Claude 3 Sonnet}: \small{The y-axis label for the subplot in row 1, column 3 is "\red{\% Win Rate}".}
\end{wronganswer}

\begin{wronganswer}
\textbf{Reka Core}: \small{ The label of the y-axis for the subplot at row 1 and column 3 is '\red{Total}'. This is explicitly written on the y-axis of the bar chart. The y-axis is shared across multiple subplots, but the label 'Total' is present on the left side of this particular subplot.}
\end{wronganswer}

\begin{wronganswer}
\textbf{InternVL Chat V1.5}: \small{The label of the y-axis for the subplot at row 1 and column 3 is "\red{\% Win Rate (c)}".}
\end{wronganswer}

\begin{wronganswer}
\textbf{MGM HD Yi 34B}: \small{The label of the y-axis for the subplot at row 1 and column 3 is '\red{Win Rate}'. This label is explicitly written and is shared across all three subplots in the first row.}
\end{wronganswer}

\begin{wronganswer}
\textbf{IDEFICS 2}: \small{\red{Recognition}.}
\end{wronganswer}
\newpage
\subsection{Example 7}
\label{subsec:desc_example7}
\begin{example}[Example]
\includegraphics[width=0.4\linewidth]{appendix_desc/533.jpg}
\begin{lstlisting}
Question: For the bottom right subplot, what is the leftmost labeled tick on the x-axis?
    * Your final answer should be the tick value on the x-axis that is explicitly written, including the case when x-axis is shared across multiple subplots. When the x-axis is present on both the top and bottom of the plot, answer based on the axis at the bottom. Ignore units or scales that are written separately from the tick, such as units and scales from the axis label or the corner of the plot.
Answer: (*$M_{\alpha^+}$*)
\end{lstlisting}
\end{example}

\begin{wronganswer}
\textbf{GPT-4o}: \small{The leftmost labeled tick on the x-axis of the bottom right subplot (D) is \red{-200}.}
\end{wronganswer}

\begin{wronganswer}
\textbf{Claude 3 Sonnet}: \small{The leftmost labeled tick on the x-axis for the bottom right subplot is \red{-Mb}.}
\end{wronganswer}

\begin{wronganswer}
\textbf{Reka Core}: \small{ The leftmost labeled tick on the x-axis of the bottom right subplot is \red{-200}. This can be determined by looking at the x-axis of the plot, where the ticks are \red{labeled with numerical values}. The tick in question is the \red{first one on the left side of the plot, which is clearly marked with the number -200}.}
\end{wronganswer}

\begin{wronganswer}
\textbf{InternVL Chat V1.5}: \small{The leftmost labeled tick on the x-axis of the bottom right subplot is \red{-200}.}
\end{wronganswer}

\begin{wronganswer}
\textbf{MGM HD Yi 34B}: \small{The leftmost labeled tick on the x-axis of the bottom right subplot is "\red{-50}". This is determined by looking at the \red{numerical values written on the x-axis}, which are aligned with the corresponding \red{bars}. The \red{"-50" tick is the first one on the left side of the x-axis}, indicating the lowest value on the scale presented in this subplot.}
\end{wronganswer}

\begin{wronganswer}
\textbf{IDEFICS 2}: \small{\red{BIC}.}
\end{wronganswer}
\newpage
\subsection{Example 8}
\label{subsec:desc_example8}
\begin{example}[Example]
\includegraphics[width=0.5\linewidth]{appendix_desc/113.jpg}
\begin{lstlisting}
Question: For the subplot at row 1 and column 5, what is the rightmost labeled tick on the x-axis?
    * Your final answer should be the tick value on the x-axis that is explicitly written, including the case when x-axis is shared across multiple subplots. When the x-axis is present on both the top and bottom of the plot, answer based on the axis at the bottom. Ignore units or scales that are written separately from the tick, such as units and scales from the axis label or the corner of the plot.
Answer: Concept
\end{lstlisting}
\end{example}

\begin{wronganswer}
\textbf{GPT-4o}: \small{The rightmost labeled tick on the x-axis for the subplot at row 1 and column 5 is "\red{Setref} Concept".}
\end{wronganswer}

\begin{wronganswer}
\textbf{Claude 3 Sonnet}: \small{The rightmost labeled tick on the x-axis for the subplot in row 1, column 5 is \red{Setref(}Concept\red{)}.}
\end{wronganswer}

\begin{wronganswer}
\textbf{Reka Core}: \small{ The rightmost labeled tick on the x-axis for the subplot at row 1 and column 5 is \red{0.15}. This is the value explicitly written on the x-axis, which represents the p-value for the Edit metric in the ShapeWorld dataset. The p-value is a statistical measure that helps determine the significance of the results, with lower values typically indicating a more statistically significant difference.}
\end{wronganswer}

\begin{wronganswer}
\textbf{InternVL Chat V1.5}: \small{The rightmost labeled tick on the x-axis for the subplot at row 1 and column 5 is \red{0.05}.}
\end{wronganswer}

\begin{wronganswer}
\textbf{MGM HD Yi 34B}: \small{The rightmost labeled tick on the x-axis of the subplot at row 1 and column 5 is "\red{0.0}". This is determined by looking at the bottom axis of the plot, where the x-axis is shared across all subplots. The tick values are written directly on the axis, and the \red{rightmost one is clearly marked as "0.0"}.}
\end{wronganswer}

\begin{wronganswer}
\textbf{IDEFICS 2}: \small{\red{Ref} Concept.}
\end{wronganswer}
\newpage
\subsection{Example 9}
\label{subsec:desc_example9}
\begin{example}[Example]
\includegraphics[width=0.5\linewidth]{appendix_desc/1860.jpg}
\begin{lstlisting}
Question: For the current plot, what is the spatially lowest labeled tick on the y-axis?
    * Your final answer should be the tick value on the y-axis that is explicitly written, including the case when y-axis is shared across multiple subplots. When the y-axis is present on both the left and right of the plot, based on the axis at the left. Ignore units or scales that are written separately from the tick, such as units and scales from the axis label or the corner of the plot.
Answer: CC daily 4.00PM
\end{lstlisting}
\end{example}

\begin{wronganswer}
\textbf{GPT-4o}: \small{The spatially lowest labeled tick on the y-axis is \red{0.2}.}
\end{wronganswer}

\begin{wronganswer}
\textbf{Claude 3 Sonnet}: \small{The spatially lowest labeled tick value on the y-axis is \red{-1.00}.}
\end{wronganswer}

\begin{wronganswer}
\textbf{Reka Core}: \small{ The spatially lowest labeled tick on the y-axis of the current plot is "\red{Alex} daily 4.00PM". This is determined by visually scanning the y-axis from top to bottom and identifying the lowest label that is clearly written and associated with the y-axis. The label "\red{Alex} daily 4.00PM" is at the bottom of the y-axis, indicating it is the lowest labeled tick.}
\end{wronganswer}

\begin{wronganswer}
\textbf{InternVL Chat V1.5}: \small{\red{-1.00}}
\end{wronganswer}

\begin{wronganswer}
\textbf{MGM HD Yi 34B}: \small{The spatially lowest labeled tick on the y-axis is \red{-0.50}. This is determined by looking at the y-axis on the left side of the plot, where the ticks are clearly marked. The lowest tick is the first one on the left, which is labeled with the value\red{-0.50}.}
\end{wronganswer}

\begin{wronganswer}
\textbf{IDEFICS 2}: \small{CC daily 4\red{\Large{\textbf{:}}}00PM.}
\end{wronganswer}
\newpage
\subsection{Example 10}
\label{subsec:desc_example10}
\begin{example}[Example]
\includegraphics[width=0.5\linewidth]{appendix_desc/1449.jpg}
\begin{lstlisting}
Question: For the subplot at row 2 and column 1, what is the spatially highest labeled tick on the y-axis?
    * Your final answer should be the tick value on the y-axis that is explicitly written, including the case when y-axis is shared across multiple subplots. When the y-axis is present on both the left and right of the plot, based on the axis at the left. Ignore units or scales that are written separately from the tick, such as units and scales from the axis label or the corner of the plot.
Answer: 0.05
\end{lstlisting}
\end{example}

\begin{wronganswer}
\textbf{GPT-4o}: \small{The spatially highest labeled tick on the y-axis for the subplot at row 2 and column 1 is \red{0.1}.}
\end{wronganswer}

\begin{wronganswer}
\textbf{Claude 3 Sonnet}: \small{The spatially highest labeled tick value on the y-axis for the subplot at row 2 and column 1 is \red{4}.}
\end{wronganswer}

\begin{wronganswer}
\textbf{Reka Core}: \small{ The spatially highest labeled tick on the y-axis for the subplot at row 2 and column 1 is \red{5}. This is determined by looking at the y-axis on the left side of the subplot, where the ticks are labeled. The tick values increase from the bottom to the top of the subplot, and the highest labeled tick is clearly marked as \red{5}.}
\end{wronganswer}

\begin{wronganswer}
\textbf{InternVL Chat V1.5}: \small{The spatially highest labeled tick on the y-axis for the subplot at row 2 and column 1 is \red{100}.}
\end{wronganswer}

\begin{wronganswer}
\textbf{MGM HD Yi 34B}: \small{The spatially highest labeled tick on the y-axis of the subplot at row 2 and column 1 is \red{2.85 m}. This is determined by looking at the y-axis on the left side of the plot, where the highest labeled tick is clearly marked.}
\end{wronganswer}

\begin{wronganswer}
\textbf{IDEFICS 2}: \small{\red{1.0}.}
\end{wronganswer}
\newpage
\subsection{Example 11}
\label{subsec:desc_example11}
\begin{example}[Example]
\includegraphics[width=0.25\linewidth]{appendix_desc/1948.jpg}
\begin{lstlisting}
Question: For the subplot at row 4 and column 2, what is difference between consecutive numerical tick values on the x-axis?
    * Your final answer should be the difference between consecutive numerical tick values of the x-axis, including the case when x-axis is shared across multiple subplots. When the x-axis is present on both the top and bottom of the plot, answer based on the axis at the bottom. Ignore units or scales that are written separately from the tick, such as units and scales from the axis label or the corner of the plot.
    * If the plot does not have an explicit x-axis tick value, or if the tick values are not numerical, or if the difference is not constant between all consecutive tick values, answer "Not Applicable".
Answer: 0.001
\end{lstlisting}
\end{example}

\begin{wronganswer}
\textbf{GPT-4o}: \small{The subplot at row 4 and column 2 has the following x-axis tick values: \red{0.00, 0.01, 0.02, 0.03, 0.04, 0.05}.

The difference between consecutive numerical tick values on the x-axis is:
\red{0.01 - 0.00 = 0.01}

Therefore, the difference between consecutive numerical tick values on the x-axis is\red{0.01}.}
\end{wronganswer}

\begin{wronganswer}
\textbf{Claude 3 Sonnet}: \small{The difference between consecutive numerical tick values on the x-axis for the subplot at row 4, column 2 is \red{0.1}.}
\end{wronganswer}

\begin{wronganswer}
\textbf{Reka Core}: \small{ The x-axis of the subplot at row 4 and column 2 is shared with the subplot to its right and has numerical tick values ranging from \red{0.00 to 0.08}. The difference between consecutive numerical tick values is consistent and can be calculated by subtracting one tick value from the next. For example, the difference \red{between 0.00 and 0.02} is 0.02, and the difference \red{between 0.02 and 0.04} is also 0.02. Therefore, the difference between consecutive numerical tick values on the x-axis is \red{0.02}.}
\end{wronganswer}

\begin{wronganswer}
\textbf{InternVL Chat V1.5}: \small{The difference between consecutive numerical tick values on the x-axis for the subplot at row 4 and column 2 is \red{0.01}.}
\end{wronganswer}

\begin{wronganswer}
\textbf{MGM HD Yi 34B}: \small{The subplot at row 4 and column 2 shows a bar chart with an x-axis that has numerical tick values. The difference between consecutive numerical tick values on the x-axis is \red{0.1}. This is determined by observing the numerical values on the x-axis, which are \red{0.0, 0.1, 0.2, 0.3, 0.4, 0.5}, and so on, indicating a constant difference of \red{0.1} between each tick.}
\end{wronganswer}

\begin{wronganswer}
\textbf{IDEFICS 2}: \small{\red{0.01}.}
\end{wronganswer}
\newpage
\subsection{Example 12}
\label{subsec:desc_example12}
\begin{example}[Example]
\includegraphics[width=0.5\linewidth]{appendix_desc/252.jpg}
\begin{lstlisting}
Question: For the subplot at row 1 and column 1, what is difference between consecutive numerical tick values on the x-axis?
    * Your final answer should be the difference between consecutive numerical tick values of the x-axis, including the case when x-axis is shared across multiple subplots. When the x-axis is present on both the top and bottom of the plot, answer based on the axis at the bottom. Ignore units or scales that are written separately from the tick, such as units and scales from the axis label or the corner of the plot.
    * If the plot does not have an explicit x-axis tick value, or if the tick values are not numerical, or if the difference is not constant between all consecutive tick values, answer "Not Applicable".
Answer: Not Applicable
\end{lstlisting}
\end{example}

\begin{wronganswer}
\textbf{GPT-4o}: \small{The difference between consecutive numerical tick values on the \red{x-axis} for the subplot at row 1 and column 1 is \red{20}.}
\end{wronganswer}

\begin{correctanswer}
\textbf{Claude 3 Sonnet}: \small{The difference between consecutive numerical tick values on the x-axis for the subplot at row 1 and column 1 is \textbf{Not Applicable}. The x-axis shows text labels ("110 items" and "agreement, 12 items") rather than numerical tick values.}
\end{correctanswer}

\begin{wronganswer}
\textbf{Reka Core}: \small{ The x-axis of the subplot at row 1 and column 1 is labeled with numerical values ranging from \red{0 to 100}. The tick marks are evenly spaced, and the difference between consecutive numerical tick values is consistent. To determine this difference, we can look at any two consecutive ticks and calculate the interval between them. For example, between the tick at 0 and the tick at 20, the difference is 20. Similarly, between the tick at 20 and the tick at 40, the difference is also 20. This pattern continues across the \red{x-axis}, indicating that the difference between consecutive numerical tick values is \red{20}.}
\end{wronganswer}

\begin{wronganswer}
\textbf{InternVL Chat V1.5}: \small{\red{10}}
\end{wronganswer}

\begin{wronganswer}
\textbf{MGM HD Yi 34B}: \small{The subplot at row 1 and column 1 has an x-axis with numerical tick values. The difference between consecutive numerical tick values on the x-axis is \red{10}. This is determined by observing the numerical labels on the \red{x-axis}, which are evenly spaced and increase by \red{10 units} between each label.}
\end{wronganswer}

\begin{wronganswer}
\textbf{IDEFICS 2}: \small{\red{1}.}
\end{wronganswer}
\newpage
\subsection{Example 13}
\label{subsec:desc_example13}
\begin{example}[Example]
\includegraphics[width=0.4\linewidth]{appendix_desc/991.jpg}
\begin{lstlisting}
Question: For the subplot at row 2 and column 1, what is difference between consecutive numerical tick values on the y-axis?
    * Your final answer should be the difference between consecutive numerical tick values of the y-axis, including the case when y-axis is shared across multiple subplots. When the y-axis is present on both the left and right of the plot, answer based on the axis at the left. Ignore units or scales that are written separately from the tick, such as units and scales from the axis label or the corner of the plot.
    * If the plot does not have an explicit y-axis tick value, or if the tick values are not numerical, or if the difference is not constant between all consecutive tick values, answer "Not Applicable".
Answer: 5
\end{lstlisting}
\end{example}

\begin{wronganswer}
\textbf{GPT-4o}: \small{The subplot at row 2 and column 1 has the y-axis tick values of -5, \red{-2.5}, 0, \red{2.5}, and 5. The difference between consecutive numerical tick values on the y-axis is \red{2.5}.}
\end{wronganswer}

\begin{wronganswer}
\textbf{Claude 3 Sonnet}: \small{The difference between consecutive numerical tick values on the y-axis for the subplot at row 2 and column 1 is \red{2.5}.}
\end{wronganswer}

\begin{wronganswer}
\textbf{Reka Core}: \small{ The y-axis tick values on the subplot at row 2 and column 1 are numerical and evenly spaced. The values are -5, \red{-4, -3, -2, -1}, 0, \red{1, 2, 3, 4}, 5. The difference between consecutive numerical tick values is \red{1}.}
\end{wronganswer}

\begin{wronganswer}
\textbf{InternVL Chat V1.5}: \small{The difference between consecutive numerical tick values on the y-axis for the subplot at row 2 and column 1 is \red{0.2}.}
\end{wronganswer}

\begin{wronganswer}
\textbf{MGM HD Yi 34B}: \small{The subplot at row 2 and column 1 shows a numerical y-axis with tick values. The difference between consecutive numerical tick values on the y-axis is 1. This is determined by observing the numerical values on the y-axis, which are -5, \red{-4, -3, -2, -1}, 0, \red{1, 2, 3, 4}, 5, and so on. The difference between any two consecutive values is \red{1}.}
\end{wronganswer}

\begin{wronganswer}
\textbf{IDEFICS 2}: \small{\red{0.1}.}
\end{wronganswer}
\newpage
\subsection{Example 14}
\label{subsec:desc_example14}
\begin{example}[Example]
\includegraphics[width=0.5\linewidth]{appendix_desc/1177.jpg}
\begin{lstlisting}
Question: For the subplot at row 1 and column 1, what is difference between consecutive numerical tick values on the y-axis?
    * Your final answer should be the difference between consecutive numerical tick values of the y-axis, including the case when y-axis is shared across multiple subplots. When the y-axis is present on both the left and right of the plot, answer based on the axis at the left. Ignore units or scales that are written separately from the tick, such as units and scales from the axis label or the corner of the plot.
    * If the plot does not have an explicit y-axis tick value, or if the tick values are not numerical, or if the difference is not constant between all consecutive tick values, answer "Not Applicable".
Answer: Not Applicable
\end{lstlisting}
\end{example}

\begin{wronganswer}
\textbf{GPT-4o}: \small{The difference between consecutive numerical tick values on the y-axis for the subplot at row 1 and column 1 is \red{0.08}.}
\end{wronganswer}

\begin{wronganswer}
\textbf{Claude 3 Sonnet}: \small{The difference between consecutive numerical tick values on the y-axis for the subplot in row 1, column 1 is \red{0.09}.}
\end{wronganswer}

\begin{wronganswer}
\textbf{Reka Core}: \small{ The y-axis of the subplot at row 1 and column 1 has numerical tick values at \red{0.04}, 0.08, \red{0.12}, and \red{0.16}. The difference between consecutive numerical tick values is \red{0.04}.}
\end{wronganswer}

\begin{wronganswer}
\textbf{InternVL Chat V1.5}: \small{The difference between consecutive numerical tick values on the y-axis of subplot at row 1 and column 1 is \red{0.01}.}
\end{wronganswer}

\begin{wronganswer}
\textbf{MGM HD Yi 34B}: \small{The difference between consecutive numerical tick values on the y-axis of the subplot at row 1 and column 1 is \red{0.08}. This is determined by observing the numerical values on the y-axis and noting the difference between the first two visible values, which are 0.08 and 0.17. Since the values are numerical and \red{the difference is constant between all consecutive tick values}, the answer is \red{not} "Not Applicable".}
\end{wronganswer}

\begin{wronganswer}
\textbf{IDEFICS 2}: \small{\red{0.05}.}
\end{wronganswer}
\newpage
\subsection{Example 15}
\label{subsec:desc_example15}
\begin{example}[Example]
\includegraphics[width=0.5\linewidth]{appendix_desc/1235.jpg}
\begin{lstlisting}
Question: For the subplot at row 1 and column 2, how many lines are there?
    * Your final answer should be the number of lines in the plot. Ignore grid lines, tick marks, and any vertical or horizontal auxiliary lines.
    * If the plot does not contain any lines or is not considered a line plot, answer "Not Applicable".
Answer: 8
\end{lstlisting}
\end{example}

\begin{wronganswer}
\textbf{GPT-4o}: \small{The subplot at row 1 and column 2 contains \red{5} lines.}
\end{wronganswer}

\begin{wronganswer}
\textbf{Claude 3 Sonnet}: \small{For the subplot in the second column, there are \red{5} lines plotted, corresponding to different values of J/U = 0.18, 0.12, 0.06, 0.0, and the reference line labeled "Ref. [14]".}
\end{wronganswer}

\begin{wronganswer}
\textbf{Reka Core}: \small{ There are \red{5} lines in the subplot at row 1 and column 2. These lines correspond to the different J/U ratios as indicated by the legend: J/U=0.18, J/U=0.12, J/U=0.06, J/U=\red{0.14}, and the reference line from Ref. [14]. Each line is represented by a different color and symbol combination, and they all follow the same path across the plot, showing the variation of the quantity \( \bar{V}^{(2)}(q) \) with respect to the momentum transfer \( q \) along the high-symmetry points in the Brillouin zone.}
\end{wronganswer}

\begin{wronganswer}
\textbf{InternVL Chat V1.5}: \small{There are \red{five} lines in the plot.}
\end{wronganswer}

\begin{wronganswer}
\textbf{MGM HD Yi 34B}: \small{The subplot at row 1 and column 2 contains \red{four} lines. These lines represent different values of the J/U ratio, as indicated by the legend in the top left corner of the image. The lines are differentiated by color and are plotted against the variable q.}
\end{wronganswer}

\begin{wronganswer}
\textbf{IDEFICS 2}: \small{\red{4}.}
\end{wronganswer}
\newpage
\subsection{Example 16}
\label{subsec:desc_example16}
\begin{example}[Example]
\includegraphics[width=0.5\linewidth]{appendix_desc/332.jpg}
\begin{lstlisting}
Question: For the current plot, how many lines are there?
    * Your final answer should be the number of lines in the plot. Ignore grid lines, tick marks, and any vertical or horizontal auxiliary lines.
    * If the plot does not contain any lines or is not considered a line plot, answer "Not Applicable".
Answer: Not Applicable
\end{lstlisting}
\end{example}

\begin{correctanswer}
\textbf{GPT-4o}: \small{\textbf{Not Applicable}}
\end{correctanswer}

\begin{wronganswer}
\textbf{Claude 3 Sonnet}: \small{The plot contains \red{4 lines}, each representing a different range of percentages: 0\%-5\%, 5\%-10\%, 10\%-20\%, and >20\%.}
\end{wronganswer}

\begin{correctanswer}
\textbf{Reka Core}: \small{ The plot in the image is a bar chart, not a line plot. Therefore, it does not contain any lines in the traditional sense of a line plot. The bars are composed of stacked segments, but these are not lines. Hence, the answer is "\textbf{Not Applicable}".}
\end{correctanswer}

\begin{wronganswer}
\textbf{InternVL Chat V1.5}: \small{\red{4}}
\end{wronganswer}

\begin{wronganswer}
\textbf{MGM HD Yi 34B}: \small{The plot contains a total of \red{four lines}, each representing a different percentage range of the data. These lines are used to differentiate the data points by their respective percentage categories.}
\end{wronganswer}

\begin{wronganswer}
\textbf{IDEFICS 2}: \small{\red{10}.}
\end{wronganswer}
\newpage
\subsection{Example 17}
\label{subsec:desc_example17}
\begin{example}[Example]
\includegraphics[width=1.0\linewidth]{appendix_desc/1007.jpg}
\begin{lstlisting}
Question: For the subplot at row 1 and column 5, do any lines intersect?
    * Your final answer should be "Yes" if any lines intersect, and "No" otherwise. Ignore grid lines, tick marks, and any vertical or horizontal auxiliary lines.
    * If the plot does not contain any lines or is not considered a line plot, answer "Not Applicable".
Answer: yes
\end{lstlisting}
\end{example}

\begin{correctanswer}
\textbf{GPT-4o}: \small{\textbf{Yes}.}
\end{correctanswer}

\begin{wronganswer}
\textbf{Claude 3 Sonnet}: \small{To determine if any lines intersect in the subplot at row 1, column 5, I will carefully examine the plot lines, ignoring grid lines, tick marks, and auxiliary lines.

In this subplot, there appears to be a \red{single line} plotted in red. The line has some fluctuations but does not intersect or cross itself at any point.

Since there is \red{only one line present} and it does not intersect itself, my answer is:

\red{No}}
\end{wronganswer}

\begin{wronganswer}
\textbf{Reka Core}: \small{ \red{No}.}
\end{wronganswer}

\begin{wronganswer}
\textbf{InternVL Chat V1.5}: \small{\red{No}, the lines in the subplot at row 1 and column 5 do not intersect.}
\end{wronganswer}

\begin{wronganswer}
\textbf{MGM HD Yi 34B}: \small{In the subplot at row 1 and column 5, the lines do not intersect. The blue line is a solid line that remains \red{above} the zero line, while the orange line is a dashed line that \red{remains below the zero line throughout the plot}. \red{There is no point where the two lines cross each other. Therefore, the answer to the question is "No".}}
\end{wronganswer}

\begin{wronganswer}
\textbf{IDEFICS 2}: \small{\red{No}.}
\end{wronganswer}
\newpage
\subsection{Example 18}
\label{subsec:desc_example18}
\begin{example}[Example]
\includegraphics[width=0.4\linewidth]{appendix_desc/1618.jpg}
\begin{lstlisting}
Question: For the bottom left subplot, do any lines intersect?
    * Your final answer should be "Yes" if any lines intersect, and "No" otherwise. Ignore grid lines, tick marks, and any vertical or horizontal auxiliary lines.
    * If the plot does not contain any lines or is not considered a line plot, answer "Not Applicable".
Answer: Not Applicable
\end{lstlisting}
\end{example}

\begin{wronganswer}
\textbf{GPT-4o}: \small{\red{No}.}
\end{wronganswer}

\begin{wronganswer}
\textbf{Claude 3 Sonnet}: \small{\red{Yes, the two lines in the bottom left subplot intersect at a point around (0.5, 0) based on the scatter plot and fitted lines shown.}}
\end{wronganswer}

\begin{wronganswer}
\textbf{Reka Core}: \small{ \red{No}.}
\end{wronganswer}

\begin{wronganswer}
\textbf{InternVL Chat V1.5}: \small{\red{No}}
\end{wronganswer}

\begin{correctanswer}
\textbf{MGM HD Yi 34B}: \small{The bottom left subplot does not contain any lines; it is a histogram. Therefore, the question of whether lines intersect is \textbf{not applicable} to this subplot.}
\end{correctanswer}

\begin{wronganswer}
\textbf{IDEFICS 2}: \small{\red{No}.}
\end{wronganswer}
\newpage
\subsection{Example 19}
\label{subsec:desc_example19}
\begin{example}[Example]
\includegraphics[width=0.5\linewidth]{appendix_desc/431.jpg}
\begin{lstlisting}
Question: For the current plot, how many discrete labels are there in the legend?
    * Your final answer should account for only labels relevant to the plot in the legend, even if the legend is located outside the plot. 
    * If the plot does not have a legend or no legend is not considered relevant to this plot, answer "Not Applicable".
Answer: 9
\end{lstlisting}
\end{example}

\begin{wronganswer}
\textbf{GPT-4o}: \small{There are \red{six} discrete labels in the legend relevant to the plot. They are:

1. Gross Output - Nordhaus
2. Gross Output - Scenario 1
3. Gross Output - Scenario 2
4. Output - Nordhaus
5. Output - Scenario 1
6. Damages - Nordhaus
7. Damages - Scenario 1
8. Damages - Scenario 2}



\end{wronganswer}

\begin{wronganswer}
\textbf{Claude 3 Sonnet}: \small{The legend for this plot contains \red{6} discrete labels: "Gross Output - Nordhaus", "Gross Output - Scenario 1", "Gross Output - Scenario 2", "Output - Nordhaus", "Output - Scenario 1", and "Damages - Scenario 2".}
\end{wronganswer}

\begin{wronganswer}
\textbf{Reka Core}: \small{ There are \red{8} discrete labels in the legend of the plot. These labels are associated with the different lines on the graph, indicating various scenarios and outputs over time. The labels are as follows: "Gross Output - Nordhaus," "Output - Nordhaus," "Damages - Nordhaus," "Gross Output - Scenario 1," "Output - Scenario 1," "Damages - Scenario 1," "Gross Output - Scenario 2," and "Output - Scenario 2." \red{The "Damages - Scenario 2" is not explicitly labeled in the legend}, but it can be inferred from the context of the other labels and the corresponding line on the graph.}
\end{wronganswer}

\begin{wronganswer}
\textbf{InternVL Chat V1.5}: \small{There are \red{8} discrete labels in the legend.}
\end{wronganswer}

\begin{wronganswer}
\textbf{MGM HD Yi 34B}: \small{The legend in the image contains \red{six} discrete labels relevant to the plot: "Gross Output - Nothaus", "Output - Nothaus", "Damages - Nothaus", "Gross Output - Scenario 1", "Output - Scenario 1", and "Damages - Scenario 1". These labels correspond to the different lines and points on the graph, indicating the data series for each category.}
\end{wronganswer}

\begin{wronganswer}
\textbf{IDEFICS 2}: \small{\red{5}.}
\end{wronganswer}
\newpage
\subsection{Example 20}
\label{subsec:desc_example20}
\begin{example}[Example]
\includegraphics[width=0.4\linewidth]{appendix_desc/384.jpg}
\begin{lstlisting}
Question: For the subplot at row 3 and column 1, how many discrete labels are there in the legend?
    * Your final answer should account for only labels relevant to the plot in the legend, even if the legend is located outside the plot. 
    * If the plot does not have a legend or no legend is not considered relevant to this plot, answer "Not Applicable".
Answer: Not Applicable
\end{lstlisting}
\end{example}

\begin{wronganswer}
\textbf{GPT-4o}: \small{The subplot at row 3 and column 1 (subplot (e)) \red{has a legend with two discrete labels}. The labels are represented by the black solid line and the red dashed line in the plot.}
\end{wronganswer}

\begin{correctanswer}
\textbf{Claude 3 Sonnet}: \small{The subplot at row 3 and column 1 does not have a legend, so the answer is "\textbf{Not Applicable}".}
\end{correctanswer}

\begin{wronganswer}
\textbf{Reka Core}: \small{ There are \red{3 discrete labels in the legend} of the subplot at row 3 and column 1. The labels are \red{"Tax Changes", "Non-Tax Topics", and "Macro Variables"}. These labels correspond to the \red{different lines in the plot, each represented by a different color and line style}.}
\end{wronganswer}

\begin{wronganswer}
\textbf{InternVL Chat V1.5}: \small{The subplot at row 3 and column 1 has \red{two discrete labels in the legend}.}
\end{wronganswer}

\begin{wronganswer}
\textbf{MGM HD Yi 34B}: \small{The subplot at row 3 and column 1 does not have a legend within the plot itself. However, \red{there is a legend located outside the plot area, which is relevant to the plot}. The legend contains \red{three discrete labels: "Contemporaneous", "1 Quarter Lag", and "2 Quarter Lag"}. Therefore, the answer to the question is \red{three}.}
\end{wronganswer}

\begin{wronganswer}
\textbf{IDEFICS 2}: \small{\red{2}.}
\end{wronganswer}
\newpage
\subsection{Example 21}
\label{subsec:desc_example21}
\begin{example}[Example]
\includegraphics[width=0.5\linewidth]{appendix_desc/177.jpg}
\begin{lstlisting}
Question: For the current plot, what are the names of the labels in the legend?
    * You should write down the labels from top to bottom, then from left to right and separate the labels with commas. Your final answer should account for only labels relevant to the plot in the legend, even if the legend is located outside the plot.
    * If the plot does not have a legend or no legend is not considered relevant to this plot, answer "Not Applicable".
Answer: min, max, max+1, max+max, 2, 5, 20
\end{lstlisting}
\end{example}

\begin{wronganswer}
\textbf{GPT-4o}: \small{min, max, max+1, max+max, \red{Unknown Vals}, 2, 5, 20}
\end{wronganswer}

\begin{wronganswer}
\textbf{Claude 3 Sonnet}: \small{The labels in the legend from top to bottom, then left to right are: Setting, min, max, max+1, max+max, \red{Unknown Vals}, 2, 5, 20.}
\end{wronganswer}

\begin{wronganswer}
\textbf{Reka Core}: \small{ The labels in the legend, from top to bottom, then from left to right, are: "Setting", "min", "max", "max+1", "max+max", "\red{Unknown Vals}", "2", "5", "20". Therefore, the final answer is: Setting, min, max, max+1, max+max, \red{Unknown Vals}, 2, 5, 20.}
\end{wronganswer}

\begin{wronganswer}
\textbf{InternVL Chat V1.5}: \small{Setting, min, max, max+1, max+max, \red{Unknown Vals}, 2, 5, 20}
\end{wronganswer}

\begin{wronganswer}
\textbf{MGM HD Yi 34B}: \small{The labels in the legend, from top to bottom and from left to right, are as follows: "min", "max", "max+1", "max+max", "\red{Unknown Vals}", "2", "5", and "20".}
\end{wronganswer}

\begin{wronganswer}
\textbf{IDEFICS 2}: \small{\red{Not Applicable}.}
\end{wronganswer}
\newpage
\subsection{Example 22}
\label{subsec:desc_example22}
\begin{example}[Example]
\includegraphics[width=0.5\linewidth]{appendix_desc/1485.jpg}
\begin{lstlisting}
Question: For the subplot at row 1 and column 1, what are the names of the labels in the legend?
    * You should write down the labels from top to bottom, then from left to right and separate the labels with commas. Your final answer should account for only labels relevant to the plot in the legend, even if the legend is located outside the plot.
    * If the plot does not have a legend or no legend is not considered relevant to this plot, answer "Not Applicable".
Answer: Not Applicable
\end{lstlisting}
\end{example}

\begin{correctanswer}
\textbf{GPT-4o}: \small{\textbf{Not Applicable}}
\end{correctanswer}

\begin{correctanswer}
\textbf{Claude 3 Sonnet}: \small{The image does not contain a legend, so the answer is "\textbf{Not Applicable}".}
\end{correctanswer}

\begin{wronganswer}
\textbf{Reka Core}: \small{ The subplot at row 1 and column 1, labeled as (a), has \red{a legend with two labels}. \red{The top label corresponds to the color blue and reads "o-pol."}, which likely stands for "sigma-polarized" or "sigma polarization." \red{The bottom label corresponds to the color red and reads "Energy loss (eV),"} indicating the scale for energy loss in electron volts (eV). \red{Therefore, the labels from top to bottom are "o-pol." and "Energy loss (eV)}."}
\end{wronganswer}

\begin{wronganswer}
\textbf{InternVL Chat V1.5}: \small{\red{Energy loss (eV), Intensity (arb. units)}}
\end{wronganswer}

\begin{wronganswer}
\textbf{MGM HD Yi 34B}: \small{The labels in the legend of the subplot at row 1 and column 1 are \red{"Intensity (arb. units)"}. These labels are relevant to the plot as they correspond to the color scale used to represent the intensity of the energy loss in the scatter plot.}
\end{wronganswer}

\begin{wronganswer}
\textbf{IDEFICS 2}: \small{\red{Incident energy (eV), Intensity (arb. units)}.}
\end{wronganswer}
\newpage
\subsection{Example 23}
\label{subsec:desc_example23}
\begin{example}[Example]
\includegraphics[width=0.5\linewidth]{appendix_desc/690.jpg}
\begin{lstlisting}
Question: For the subplot at row 5 and column 2, what is the difference between the maximum and minimum values of the tick labels on the continuous legend (i.e., colorbar)?
    * You should remove the percentage sign (if any) in your answer.
    * If the plot does not have an explicit colorbar-based continuous legend or the legend is not considered relevant to this subplot, answer "Not Applicable".
Answer: 0.8
\end{lstlisting}
\end{example}

\begin{wronganswer}
\textbf{GPT-4o}: \small{The subplot at row 5 and column 2 \red{does not have an explicit colorbar-based continuous legend}. Therefore, the answer is \red{"Not Applicable"}.}
\end{wronganswer}

\begin{wronganswer}
\textbf{Claude 3 Sonnet}: \small{\red{Not Applicable}}
\end{wronganswer}

\begin{wronganswer}
\textbf{Reka Core}: \small{ The continuous legend in the subplot at row 5 and column 2, \red{labeled C2S, ranges from 0 to 200}. The maximum value is 200 and the minimum value is 0. Therefore, the difference between the maximum and minimum values of the tick labels on the continuous legend is \red{200}.}
\end{wronganswer}

\begin{wronganswer}
\textbf{InternVL Chat V1.5}: \small{The subplot at row 5 and column 2 has a continuous legend (colorbar) on the right side of the plot. \red{The tick labels on the colorbar range from 0 to 200}. The difference between the maximum and minimum values of the tick labels is 200 - 0 = \red{200}.}
\end{wronganswer}

\begin{wronganswer}
\textbf{MGM HD Yi 34B}: \small{The subplot at row 5 and column 2 \red{does not have an explicit colorbar-based continuous legend}. Instead, it has a scatter plot with a color-coded legend on the right side, which is not relevant to the question about the difference between the maximum and minimum values of the tick labels on the continuous legend. Therefore, the answer is "\red{Not Applicable}".}
\end{wronganswer}

\begin{wronganswer}
\textbf{IDEFICS 2}: \small{\red{Not Applicable}.}
\end{wronganswer}
\newpage
\subsection{Example 24}
\label{subsec:desc_example24}
\begin{example}[Example]
\includegraphics[width=0.5\linewidth]{appendix_desc/1726.jpg}
\begin{lstlisting}
Question: For the bottom right subplot, what is the difference between the maximum and minimum values of the tick labels on the continuous legend (i.e., colorbar)?
    * You should remove the percentage sign (if any) in your answer.
    * If the plot does not have an explicit colorbar-based continuous legend or the legend is not considered relevant to this subplot, answer "Not Applicable".
Answer: Not Applicable
\end{lstlisting}
\end{example}

\begin{wronganswer}
\textbf{GPT-4o}: \small{The bottom right subplot \red{has a continuous legend (colorbar) with tick labels ranging from 0.8 to 2.3}. The difference between the maximum and minimum values of the tick labels is:

2.3 - 0.8 = 1.5

So, the answer is \red{1.5}.}
\end{wronganswer}

\begin{wronganswer}
\textbf{Claude 3 Sonnet}: \small{The difference between the maximum and minimum values \red{on the continuous legend (colorbar) for the bottom right subplot is 1.5}.}
\end{wronganswer}

\begin{wronganswer}
\textbf{Reka Core}: \small{ The difference between the maximum and minimum values of the tick labels \red{on the continuous legend for the bottom right subplot is 2.3 - 0.8 = 1.5. The continuous legend, or colorbar, ranges from 0.8 to 2.3, and the values are associated with the parameter $\beta$}, which is part of the Generalised Fisher model. The colorbar is relevant to this subplot as it indicates the range of $\beta$ values corresponding to the different colors on the plot.}
\end{wronganswer}

\begin{wronganswer}
\textbf{InternVL Chat V1.5}: \small{The maximum value of the tick labels \red{on the continuous legend in the bottom right subplot is 2.3}, and the minimum value is 0. The difference between these values is 2.3 - 0 = \red{2.3}.}
\end{wronganswer}

\begin{wronganswer}
\textbf{MGM HD Yi 34B}: \small{The bottom right subplot of the "Generalised Fisher" section \red{shows a continuous legend} with tick labels ranging from 0.8 to 2.3. The maximum value is 2.3, and the minimum value is 0.8. Therefore, the difference between the maximum and minimum values of the tick labels on the continuous legend is \red{1.5}.}
\end{wronganswer}

\begin{wronganswer}
\textbf{IDEFICS 2}: \small{\red{0.1}.}
\end{wronganswer}
\newpage
\subsection{Example 25}
\label{subsec:desc_example25}
\begin{example}[Example]
\includegraphics[width=0.5\linewidth]{appendix_desc/1003.jpg}
\begin{lstlisting}
Question: For the subplot at row 2 and column 1, what is the maximum value of the tick labels on the continuous legend (i.e., colorbar)?
    * You should remove the percentage sign (if any) in your answer. 
    * If the plot does not have an explicit colorbar-based continuous legend or the legend is not considered relevant to this subplot, answer "Not Applicable".
Answer: 0.8
\end{lstlisting}
\end{example}

\begin{wronganswer}
\textbf{GPT-4o}: \small{The subplot at row 2 and column 1 is labeled "c" and is titled "Probability of static coverage." The maximum value of the tick labels on the continuous legend (colorbar) for this subplot is \red{1.0}.}
\end{wronganswer}

\begin{wronganswer}
\textbf{Claude 3 Sonnet}: \small{\red{Not Applicable}}
\end{wronganswer}

\begin{wronganswer}
\textbf{Reka Core}: \small{ \red{35}.}
\end{wronganswer}

\begin{wronganswer}
\textbf{InternVL Chat V1.5}: \small{\red{35}}
\end{wronganswer}

\begin{wronganswer}
\textbf{MGM HD Yi 34B}: \small{The maximum value of the tick labels on the continuous legend (colorbar) for the subplot at row 2 and column 1 is \red{$10^1$}. This is determined by observing the colorbar on the right side of the subplot, where the highest value is indicated by the darkest color, which corresponds to the largest number of mobile sensors.}
\end{wronganswer}

\begin{wronganswer}
\textbf{IDEFICS 2}: \small{\red{Not Applicable}.}
\end{wronganswer}
\newpage
\subsection{Example 26}
\label{subsec:desc_example26}
\begin{example}[Example]
\includegraphics[width=0.5\linewidth]{appendix_desc/1386.jpg}
\begin{lstlisting}
Question: For the subplot at row 5 and column 3, what is the maximum value of the tick labels on the continuous legend (i.e., colorbar)?
    * You should remove the percentage sign (if any) in your answer. 
    * If the plot does not have an explicit colorbar-based continuous legend or the legend is not considered relevant to this subplot, answer "Not Applicable".
Answer: Not Applicable
\end{lstlisting}
\end{example}

\begin{correctanswer}
\textbf{GPT-4o}: \small{\textbf{Not Applicable}}
\end{correctanswer}

\begin{correctanswer}
\textbf{Claude 3 Sonnet}: \small{\textbf{Not Applicable}}
\end{correctanswer}

\begin{wronganswer}
\textbf{Reka Core}: \small{ The subplot at row 5 and column 3 corresponds to the supernova labeled \red{"14C lb"}. The continuous legend in this subplot is \red{represented by the colorbar on the right side of the figure, which ranges from 0 to 1. The maximum value of the tick labels on this colorbar is 1}.}
\end{wronganswer}

\begin{wronganswer}
\textbf{InternVL Chat V1.5}: \small{The maximum value of the tick labels \red{on the continuous legend (i.e., colorbar)} for the subplot at row 5 and column 3 is \red{30}.}
\end{wronganswer}

\begin{wronganswer}
\textbf{MGM HD Yi 34B}: \small{The subplot at row 5 and column 3 is labeled \red{"09fj lb"}. \red{The continuous legend (colorbar)} for this subplot has tick labels ranging from 0.4 to 0.8. The maximum value of the tick labels on the continuous legend is \red{0.8}.}
\end{wronganswer}

\begin{wronganswer}
\textbf{IDEFICS 2}: \small{\red{1.6}.}
\end{wronganswer}
\newpage
\subsection{Example 27}
\label{subsec:desc_example27}
\begin{example}[Example]
\includegraphics[width=0.5\linewidth]{appendix_desc/1524.jpg}
\begin{lstlisting}
Question: For the subplot at row 1 and column 2, what is the general trend of data from left to right?
    * Your final answer should be within a few words, such as "increases", "increases then stabilizes".
Answer: decrease then increase
\end{lstlisting}
\end{example}

\begin{wronganswer}
\textbf{GPT-4o}: \small{\red{Increases}.}
\end{wronganswer}

\begin{wronganswer}
\textbf{Claude 3 Sonnet}: \small{The general trend of the data for men in the top right subplot is that it \red{increases sharply from left to right across the age range shown}.}
\end{wronganswer}

\begin{correctanswer}
\textbf{Reka Core}: \small{ The general trend of the data from left to right in the subplot at row 1 and column 2 is that \textbf{the log hazard rate decreases as age increases, reaches a minimum around age 50, and then increases again with age}. This is indicative of a U-shaped curve.}
\end{correctanswer}

\begin{wronganswer}
\textbf{InternVL Chat V1.5}: \small{\red{increases}}
\end{wronganswer}

\begin{wronganswer}
\textbf{MGM HD Yi 34B}: \small{The general trend of data from left to right in the subplot at row 1 and column 2 is "\red{increases}". This is observed as the hazard rate for both 1885 and 1905 data points rises with age, indicating a higher risk of death as age increases.}
\end{wronganswer}

\begin{wronganswer}
\textbf{IDEFICS 2}: \small{\red{Increases}.}
\end{wronganswer}
\newpage
\subsection{Example 28}
\label{subsec:desc_example28}
\begin{example}[Example]
\includegraphics[width=0.5\linewidth]{appendix_desc/20.jpg}
\begin{lstlisting}
Question: For the current plot, What is the total number of explicitly labeled ticks across all axes?
    * Your final answer should be the total number of explicitly labeled ticks across all axes, including the case when any axis is shared across multiple subplots.
Answer: 16
\end{lstlisting}
\end{example}

\begin{wronganswer}
\textbf{GPT-4o}: \small{To determine the total number of explicitly labeled ticks across all axes in the given plot, we need to count the labeled ticks on both the x-axis and the y-axis.

1. **X-axis (Degree):**
   - The labeled ticks are: 0, 50, 100, 150, 200, 250, 300
   - Total labeled ticks on the x-axis: 7

2. **Y-axis (Proportion of papers):**
   - \red{The labeled ticks are: 0.00, 0.10, 0.20, 0.30, 0.40}
   - Total labeled ticks on the y-axis: \red{5}

Adding these together, the total number of explicitly labeled ticks across all axes is:
\[ 7 \text{ (x-axis)} + \red{5} \text{ (y-axis)} = \red{12} \]

So, the total number of explicitly labeled ticks across all axes is \red{12}.}
\end{wronganswer}

\begin{wronganswer}
\textbf{Claude 3 Sonnet}: \small{The total number of explicitly labeled ticks across all axes in the given plot is \red{7}. \red{There is one axis} labeled "Degree" with ticks at 0, 50, 100, 150, 200, 250, and 300.}
\end{wronganswer}

\begin{wronganswer}
\textbf{Reka Core}: \small{ The total number of explicitly labeled ticks across all axes in the current plot is \red{11}. This includes the x-axis, which has 7 labeled ticks from 0 to 300 in increments of 50, and the y-axis, \red{which has 5 labeled ticks} from 0.00 to 0.40 in increments of 0.05. There are no additional subplots or shared axes in this single plot.}
\end{wronganswer}

\begin{wronganswer}
\textbf{InternVL Chat V1.5}: \small{\red{6}}
\end{wronganswer}

\begin{wronganswer}
\textbf{MGM HD Yi 34B}: \small{The total number of explicitly labeled ticks across all axes in the current plot is \red{10}. This includes \red{the 5 ticks on the x-axis and the 5 ticks on the y-axis}. There are no shared axes or subplots in this image, so the count is straightforward.}
\end{wronganswer}

\begin{wronganswer}
\textbf{IDEFICS 2}: \small{\red{3}.}
\end{wronganswer}
\newpage
\subsection{Example 29}
\label{subsec:desc_example29}
\begin{example}[Example]
\includegraphics[width=0.5\linewidth]{appendix_desc/2287.jpg}
\begin{lstlisting}
Question: What is the layout of the subplots?
    * Your final answer should follow "n by m" format, where n is the number of rows and m is the number of columns.
    * If the plot does not contain subplots, answer "1 by 1".
Answer: 5 by 1
\end{lstlisting}
\end{example}

\begin{wronganswer}
\textbf{GPT-4o}: \small{The layout of the subplots is \red{4} by 1.}
\end{wronganswer}

\begin{wronganswer}
\textbf{Claude 3 Sonnet}: \small{The layout of the subplots in the given image is \red{6} by 1, meaning there are \red{6} rows and 1 column of subplots.}
\end{wronganswer}

\begin{wronganswer}
\textbf{Reka Core}: \small{ The image contains a grid of subplots arranged in a \red{3 by 3 format, with three rows and three columns}. Each subplot is labeled with a different expectation value, and the plots are color-coded to represent two different methods, DML and C-DML. The x-axis of all subplots is labeled with $\sigma_u$, which suggests that the plots are showing the relationship between these expectation values and $\sigma_u$. The y-axes vary depending on the subplot, with different scales and units, indicating different metrics being compared.}
\end{wronganswer}

\begin{wronganswer}
\textbf{InternVL Chat V1.5}: \small{The layout of the subplots is \red{4} by 1.}
\end{wronganswer}

\begin{wronganswer}
\textbf{MGM HD Yi 34B}: \small{The layout of the subplots is \red{"2 by 3"}. This means there are \red{two rows and three columns of subplots}. Each subplot contains two lines, one red and one blue, representing different data sets or conditions.}
\end{wronganswer}

\begin{wronganswer}
\textbf{IDEFICS 2}: \small{\red{1 by 5}.}
\end{wronganswer}
\newpage
\subsection{Example 30}
\label{subsec:desc_example30}
\begin{example}[Example]
\includegraphics[width=0.5\linewidth]{appendix_desc/266.jpg}
\begin{lstlisting}
Question: What is the number of subplots?
    * Your final answer should be the total number of subplots in the plot.
    * If the plot does not contain subplots, answer "1".
Answer: 11
\end{lstlisting}
\end{example}

\begin{wronganswer}
\textbf{GPT-4o}: \small{The plot contains \red{10} subplots.}
\end{wronganswer}

\begin{wronganswer}
\textbf{Claude 3 Sonnet}: \small{The image contains \red{12 subplots arranged in a grid layout}.}
\end{wronganswer}

\begin{correctanswer}
\textbf{Reka Core}: \small{ The image contains a total of \textbf{11 subplots}. Each subplot represents a different dataset, as indicated by the labels on the top of each subplot. These datasets include ImageNet-val, Omniglot, Aircraft, Birds, Textures, Quick Draw, Fungi, VGG Flower, Traffic Signs, MSCOCO, and a combined plot for the average test error over 9 datasets. Therefore, the final answer is \textbf{11 subplots}.}
\end{correctanswer}

\begin{wronganswer}
\textbf{InternVL Chat V1.5}: \small{\red{9}}
\end{wronganswer}

\begin{wronganswer}
\textbf{MGM HD Yi 34B}: \small{The plot contains a total of \red{10} subplots. Each subplot represents a different dataset or category, as indicated by the labels in the top row of the plot. The subplots are arranged in \red{two columns and five rows, making a total of ten individual plots}.}
\end{wronganswer}

\begin{wronganswer}
\textbf{IDEFICS 2}: \small{\red{6}.}
\end{wronganswer}